\documentclass[11pt]{article}
\newif\ifdraft\drafttrue
\usepackage{times}
\usepackage{abbrevs}
\usepackage{balance}
\usepackage{alltt}
\usepackage{graphicx}
\usepackage{txfonts}
\usepackage{pgfgantt}
\usepackage{afterpage}
\usepackage{enumitem}
\usepackage{booktabs}
\usepackage{float}
\usepackage{natbib}
\setlength{\bibsep}{0.0pt}


\newif\ifcolor\colortrue

% Color
\usepackage{xcolor}
\definecolor{pennblue}{cmyk}{1,.65,0,.3}
\definecolor{pennred}{cmyk}{0,1,.65,.34}
\definecolor{Gray}{gray}{0.9}

\usepackage[margin=1in,includefoot]{geometry}

% Hyperlinks
\ifcolor
\usepackage[colorlinks=true,linkcolor={pennblue},citecolor={pennred},urlcolor={pennblue}]{hyperref}
\else
\usepackage[colorlinks=true,linkcolor={black},citecolor={black},urlcolor={black}]{hyperref}
\fi
\usepackage{url}


\definecolor{red3}{rgb}{0.80,0.00,0.00}
\newcommand{\note}[1]{[\textcolor{red}{\textit{#1}}]}
\newcommand{\rjs}[1]{\note{RS: #1}}
\newcommand{\TODO}[1]{\textcolor{red3}{[{\bfseries TODO: }{#1}]}}



\def\PAR#1{\par\medskip\noindent\textbf{#1}}

\begin{document}

%% Your proposal should outline the specific problem and goals of your research, and ideally, be limited to 2 pages.


%% - A description of the planned testing environment and how it will demonstrate advances in networking capabilities relevant to the research context.
%% - A plan for conducting periodic demonstrations.
%% - A description of metrics to be used to assess the performance of their systems, along with a timetable and methodology for such assessments.

\section*{Intel Connectivity: A Hardware Accelerator for Key-Value (KV) Caches}

\PAR{Principal Investigators:} Robert Soul{\'e} and Rajit Manohar
(\{robert.soule, rajit.manohar\}@yale.edu), Departments of Computer Science and Computer Engineering,
Yale University 10 Hillhouse Avenue,, New Haven, CT, 06511.

\PAR{Intel contact and sponsor:} James Tsai
  (\href{james.tsai@intel.com}{james.tsai@intel.com})


%% At the top of the hierarchy are
%% SRAM caches and DRAM main memory, which have low latency, unlimited
%% write endurance, and fine granularity of access.  They are, however,
%% power-hungry, expensive and volatile, necessitating further tiering to
%% solid-state NAND flash storage (SSD), and finally to spinning disk or
%% tape magnetic storage at the bottom of the hierarchy. These terminal
%% tiers of non-volatile and durable bit storage have much higher access
%% latency and granularity than volatile memory.



\PAR{Introduction.}
%
There is intense competition to build the lowest-latency inference pipelines.
Companies such as Cerebras, Grok, Groq, and Etched have developed specialized
hardware accelerators that threaten Nvidia's GPU hegemony on compute platforms
and ...

As a result, the current performance bottleneck for the decode phase is not
compute, but rather, memory bandwidth (limited by loading model weights and the growing cache from memory).




Attention~\cite{vaswani2017}


\PAR{Problem Statement.}
%
The central problem motivating this proposal is: \emph{what
  will be needed from network protocols and services to support
large pools of disaggregated storage class memory?}


\PAR{Research Goals.}
%
Expressed succinctly, the research goals of this proposal are to:
\begin{enumerate}[nolistsep]
\item ...
\item ...
\item ...
\item ...

\end{enumerate}



\PAR{Research Team.}
%
The proposed work will be carried out by a team
of three researchers: the two faculty PIs, and one PhD
student. PI Soul\'{e} is a
recognized expert in data plane programming.  He spent a sabbatical at
Barefoot Networks, and has published extensively on the topic.
Professor Manohar is an expert in asynchronous hardware
design. 


both operating systems, databases,
and distributed systems.
PI Silberschatz is a fellow of the ACM, a fellow the IEEE, and a member
of the Connecticut Academy of Science and Engineering.
He is the author of two well-known textbooks, ``Operating
System Concepts'' and “Database System Concepts.”



\PAR{Test Environment.}
%%
The PIs maintain a small cluster of Dell R740 servers with 1.5TB of
Optane memory. To experiment and evaluate new network protocols and
support, we will rely eBPF to modify the kernel, or a set of Pensando
25G NICs that have been donated to us. We will need Cisco Tofino-based
Nexus switches to evaluate different protocols in the network core. We
will therefore need to acquire 2x Cisco Nexus 34180YC Programmable
switches.


\PAR{Collaboration Plan and Progress.}
%%
To demonstrate our progress, we will have conference calls with our Intel Sponsor, James Tsai, on a bi-weekly basis. 

%% To avoid semantic conflict between user vs. cisco p4 programs, some of the metadata should be hidden or made immutable for the user pipe, but I guess P4_16 already can handle such needs.

\PAR{Expected Outcome.}
%
After one year of work, we expect to deliver: an operating
system kernel modification that can issue remote memory requests that include the
 permanent pointer. The second milestone will be to route
 these memory requests across multiple hops (i.e., more than
 one forwarding element). The third milestone will be to map
 these requests to ABD consensus requests across replicated
 memory instances. A forth (and stretch) goal will be to add
 support for capabilities.


%% \PAR{Data Policy.} We will publicly release as open-source all code,
%% as well as the data sets used to evaluate our system, except for those
%% components and data that would reveal proprietary information from our
%% partners.

{\tiny
\bibliographystyle{abbrv}
\balance
\bibliography{main}
}

\end{document}

